\sectionguard
\section{Scope, Status, and Qualifications}

\subsection*{Question and thesis boundary}

This study answers one question: what was reported at Lourdes in
1858 on the nearest genuinely accessible record; what exactly the
competent ecclesiastical authority judged on 18 January 1862, in
which words and with what object; how the reported cures are and
have been examined and recognized; how Bernadette, the liturgy,
and the popes received the event; and what all of that obliges
and leaves open. It is a source-audited study instrument,
ecclesial-judgment-first in method: not a devotional narrative, a
critical edition, a magisterial act, a medical investigation, or
a complete history of Lourdes. It proposes no judgment of its own
on supernaturality and re-diagnoses no medical file.

\subsection*{Corpus and exclusions}

Included: the eighteen reported apparitions of 11 February--16
July 1858, bounded to the witness strata named in \S2.1; the
civil interrogations at their documented existence; the
1858--1862 commission and the mandement, quoted from Lasserre's
1870 reproduction collated with the Sanctuary's transcriptions;
the self-identification's textual history at the accessible
witnesses; the medical-recognition institutions and the
Sanctuary's published list of 72 recognized cures with the
dossiers of the 70th--72nd; Bernadette's later life through the
1933 canonization acts in AAS 26 (1934); liturgical reception
(1890, 1907, 1969-era rank at ceiling, 1992); papal reception
(Pius IX to Benedict XVI at the cited texts; the announced 2026
visit as announcement only); and current standing under the 2024
DDF norms. Excluded or bounded: any certification of physical or
medical mechanism; any reconstruction of the personal prayer or
secrets; the candle episode beyond received narrative; the
contents of the police and judicial dossiers (not consulted);
Estrade and the memoir literature (not consulted); the
nineteenth-century controversies (named as reception history
only); apparitions claimed elsewhere; and every claim of the
comprehensive ``all approved apparitions'' type---this is a
single-event study and asserts no registry.

\subsection*{Geography, chronology, and currentness}

Subject geography: Lourdes and the grotto of Massabielle;
Bartr\`es; Tarbes; Nevers; Rome; the dioceses of the recognized
cures. Subject range: 1844 (Bernadette's birth) to the acts and
announcements current at the as-of date. As-of date: all mutable
claims---current controlling status, the DDF documents index,
Sanctuary pages, the cure count of 72, the Bureau staffing, the
announced papal visit---were checked on 2026-07-25. No later act
superseding the 1862 judgment was located. Stable historical
facts are not given artificial currentness. Negative results are
bounded searches, correctable by better evidence, never
guarantees. Two Sanctuary-site internal stalenesses found during
research are recorded rather than harmonized: a general page
still naming Dr de Franciscis as currently responsible after the
dated announcement of Dr Monami's appointment, and an older
statistics page still counting 70 cures.

\subsection*{Taxonomy and terminology}

Kept distinct throughout: \emph{received custodial narrative}
(Sanctuary chronology and pages); \emph{period narrative
literature} (Lasserre); \emph{documentary reproduction} (the
mandement as printed by Lasserre; the AAS/ASS acts);
\emph{competent judgment} (the 1862 act; each cure recognition;
the canonization); \emph{medical finding} (``unexplained in the
present state of medical knowledge''---never more);
\emph{liturgical reception} (feast, World Day of the Sick);
\emph{papal teaching and act of piety} (encyclical, homilies,
visits); \emph{announcement} (the 2026 visit); and
\emph{editorial synthesis}. ``Reported'' is the default register
for all words and phenomena attributed to the apparition.
\term{Fond\'es \`a la croire certaine} is the 1862 act's own
formula and is not translated into the 2024 taxonomy.
\term{Fides humana} names the prudent human assent invited by
approved private revelation.

\subsection*{Source hierarchy and method}

Priority: the exact operative and documentary texts in
identified witnesses (the mandement in the 1870 reproduction
with page-image collation; AAS 26 (1934); ASS 40 (1907); the
Holy See's web texts for 1957--2024); then the Sanctuary's and
AMIL's institutional pages for the custodial narrative, the
medical institutions, and the cures list, always at
institutional-self-description ceilings; then period narrative
(Lasserre) for what only it witnesses, marked as such; then
common public record, marked as such and used sparingly.
Quotations are exact from the named witness; French, Latin,
Italian, and Occitan sources are quoted in the original where
wording is operative, with editorial translations visibly
editorial. No project or AI rendering is offered as a received
translation; no prayer text is printed for recitation in this
study.

\subsection*{Qualifications and unresolved items}

(1)~The 1862 mandement was read in Lasserre's 1870 reproduction
(pages collated as images) and the Sanctuary's transcriptions,
not in the diocesan original or 1862 imprint; Lasserre abridges
the narrative portion and arts.~3--4. (2)~The critical edition
of the 1858 documents (Laurentin--Billet, \work{Documents
authentiques}) and Bernadette's autograph \'ecrits were not
accessible; the police and judicial records are asserted only as
to documented existence, and officials' forenames are not
asserted. (3)~Estrade's memoir was not consulted. (4)~The Gascon
sentence's orthography varies between the Sanctuary's own pages;
no autograph spelling is asserted. (5)~The apparitions table
rests on the Sanctuary's chronology, a custodial harmonization;
crowd figures are received round numbers. (6)~AAS 26 (1934) and
ASS 40 (1907) were read in the Holy See's OCR scans; OCR defects
are noted where material (the canonization formula's
\term{exaltationem}). (7)~The current rank of the 11 February
memorial and the French particular date for St Bernadette are
reported at ceiling without typical-edition verification.
(8)~The cures list is the Sanctuary's; its ordinals are its
order anchored by its own 70th--72nd designations; cure dates
are asserted only where a dedicated institutional article gives
them; the aggregate statistics describe the 70-cure state.
(9)~The 1905 Roman act confirming the Bureau's procedures and
the pre-1890 privilege acts were not located as documents; the
Sanctuary's summaries are used at ceiling. (10)~The 2026 papal
visit is an announcement and may change. (11)~No independent
Mariological, historical, canonical, archival, Occitan or French
philological, liturgical, medical, or scientific review has
occurred; no ecclesiastical review or approval is claimed or
implied.

\subsection*{Rights}

Project prose is original and intended for CC~BY~4.0. Quotations
remain under their own status: the 1862 mandement and the other
nineteenth-century French and Latin documents are public domain
in text (Lasserre died in 1900; official acts), quoted with
citation; \work{Acta Apostolicae Sedis} and \work{Acta Sanctae
Sedis} material and Holy See web texts are quoted within
ordinary scholarly limits with citation; Sanctuary of Lourdes
and AMIL institutional pages are \copyright\ their publishers
and are quoted briefly with citation; the Papal Encyclicals
Online rendering of \work{Ineffabilis Deus} has an unidentified
translator and unresolved translation rights and is quoted once,
briefly, with attribution. Online availability is not treated as
a reuse license; no complete document is reproduced; scan images
were consulted for collation and are not redistributed.
